\documentclass{article}
\usepackage{tikz}
\usetikzlibrary{angles, quotes}
\usepackage{amssymb}
\usepackage{pgfplots}
\pgfplotsset{width=10cm,compat=1.9}
\usepgfplotslibrary{statistics}
\usepackage[left=2.54cm,right=2.54cm,top=2.54cm,bottom=2.54cm]{geometry}
\usepackage{tasks} %Paquete necesario para la producción de items horizontales para algunos ejercicios de tarea empleando el comando /begin{multicols}}{#}.../end{multicols} para ayudarnos a reducir las páginas
\usepackage{pdflscape} %comando que permite cambiar a orientación horizontal las hojas%
\usepackage{soul} %Paquete para permitir la compilación de acentos usando el comando {\hl{}}}\\
\sethlcolor{yellow(munsell)} %comando para definir el color para el subrayado usando el comando \hl
\usepackage[utf8]{inputenc}
\usepackage[spanish]{babel}
\usepackage{icomma}
\usepackage{ragged2e}		
\usepackage{siunitx}
\usepackage{url}
\usepackage[colorlinks=true, urlcolor=blue,  linkcolor=Cerulean, citecolor=green]{hyperref}
\usepackage{pdfpages}
\usepackage{blindtext} %Paquete que genera texto ficticio (dummy text) con el comando: \blindtext
\usepackage{cancel} %in the preamble gives you four different modes of striking through: \cancel{text to cancel} draws a diagonal line (slash) through its argument, \bcancel{text to cancel} uses the negative slope (a backslash), \xcancel{text to cancel} draws an X (actually \cancel plus \bcancel), \cancelto{〈value〉}{〈expression〉} draws a diagonal arrow through the 〈expression〉pointing to the 〈value〉 (math-mode only)
\usepackage{csquotes}
\usepackage{afterpage}
\usepackage{parskip} 
\usepackage{float}
\usepackage{enumitem}
%\usepackage{multicol}%Paquete que permite la creación aislada de columnas de texto. De esta forma se reduce la cantidad de páginas en nuestro documento
\usepackage{lipsum} %Paquete que genera texto de relleno al igual que \usepackage{blindtext}, se usa con el comando \lipsum[#-#] los #(hashtags) delimitan el número de páginas que deseamos ocupar del paquete lipsum para usarlos como relleno. 

% Definir el comando personalizado para la nota
\newcommand{\mynote}[1]{%
\begin{tikzpicture}[baseline=-0.75ex]
\node[draw, circle, fill=Apple Green!60, inner sep=2pt] (note) {\textbf{Nota:}};
\end{tikzpicture}%
\ \textit{#1}%
}

\newenvironment{Figura}
{\par\medskip\noindent\minipage{\linewidth}}
{\endminipage\par\medskip}
\usepackage{caption}
\usepackage[
backend=biber,
sorting=none,
url=true
]{biblatex} %bibliografía
\addbibresource{biblio.bib}
% Establecer el espacio entre las entradas de la bibliografía
\setlength{\bibitemsep}{1\baselineskip} % Puedes ajustar el valor según tus preferencias
\usepackage{amsmath,amsthm,amssymb,amsfonts}
\usepackage{pifont} %Permite la compilación del comando \xmark: es el símbolo de la tachita
\usepackage{mathtools} %permite la compilación de símbolos para matrices
\usepackage{empheq} %Paquete que se relaciona con \usepackage[most]{tcolorbox} para la creación de cuadros/cajas de colores para delimitar los resultados matemáticos o líneas de texto usando la línea de comando: \begin{empheq}[box={\mymath[colback=orange(webcolor)!70,drop lifted shadow]}]{equation*} \end{empheq}
\usepackage[most]{tcolorbox} %Paquete que se relaciona con \usepackage{empheq} para la creación de cuadros/cajas de colores para delimitar los resultados matemáticos o líneas de texto
\renewcommand{\qedsymbol}{$\blacksquare$}
\usetikzlibrary{positioning,decorations.pathreplacing} %Paquete que permite la compilación de llaves para cuadros sinópticos (brace diagram) 
\usepackage{schemata} %The schemata package is designed just for these brace diagrams
\usepackage{booktabs}

\newcommand\AB[2]{\schema{\schemabox{#1}}{\schemabox{#2}}} %comando o paquete necesario para crear cuadros sinópticos

\newtcbox{\mymath}[1][]{%
nobeforeafter, math upper, tcbox raise base,
enhanced, colframe=blue!30!black,
colback=blue!30, boxrule=1pt,
#1} %Comando importante para encasillar los resultados matemáticos en cajas de diferentes colores, se relaciona con los paquetes \usepackage{empheq} y \usepackage[most]{tcolorbox} 

\newenvironment{sysmatrix}[1]
{\left(\begin{array}{@{}#1@{}}}
{\end{array}\right)}
\newcommand{\ro}[1]{%
\xrightarrow{\mathmakebox[\rowidth]{#1}}%
}
\newlength{\rowidth}% row operation width
\AtBeginDocument{\setlength{\rowidth}{3em}}  %Comando importante para laproducción de lineas de operaciones entre matrices, método de Gauss o Gauss Jordan se relaciona con \newenvironment{sysmatrix}

%formato para cambiar el horario a español
\usepackage[spanish]{datetime2}
\DTMsetdatestyle{spanish}

\renewcommand{\today}{\DTMdisplaydate{\the\year}{\the\month}{\the\day }{-1}}
%%%%%%%%%%%%%%%%%%%%%%%%%%%%%%%%%%%%%%%%%

\usepackage{datetime}
\newcommand{\mycurrenttime}{\xxivtime}

%%%%%%%%%%%%%%%%%%%%%%%%%%%%%%%%%%%%%%%%%%%%%%%%%%%%%%%%%%%%%%%%%%%%%%%%%%%%%%%%%%%%%%%%%%%%%%%%%%%%%%%%%%%%%%%%%%%%%%%%%%%%%%%%%%%%%%%%%%%%%%%%%


%%%%%%%%%%%%%Caja de color para el título%%%%%%%%%%%%%%%%%%%%%%%%%%%%%%%%%%%%%%%%%%%%%%%%%%%%%%

\definecolor{myframecolor}{RGB}{1, 45, 75} % Prussian Blue
\definecolor{myboxcolor}{RGB}{0, 107, 92} % Apple Green

\newtcolorbox{mybox}{
enhanced,
colback=myboxcolor!7, % Color de fondo del cuadro
colframe=myframecolor, % Color del marco del cuadro
arc=0pt,
boxrule=1pt,
borderline west={2mm}{-10mm}{myframecolor}, % Borde en el lado izquierdo
sharp corners=southwest,
width=\linewidth,
before=\par\vspace{\bigskipamount}, % Espacio antes del cuadro
after=\par\vspace{\bigskipamount} % Espacio después del cuadro
}

%%%%%%%%%%%%%%%%%%%%%%%%%%%%%%%%%%%%%%%%%%%%%%%%%%%%%%%%%Nueva caja%%%%%%%%%%%%%%%%%%%%%%%%%%%%%%%%%%%%%%%%%%%%%%%%%%%%%%%%%%%%%%%%%

\newtcolorbox{mybox2}{
enhanced,
colback=Apple Green!7, % Color de fondo del cuadro
colframe=Apple Green, % Color del marco del cuadro
arc=0pt,
boxrule=1pt,
borderline west={2mm}{-10mm}{Apple Green}, % Borde en el lado izquierdo
sharp corners=southwest,
width=\linewidth,
before=\par\vspace{\bigskipamount}, % Espacio antes del cuadro
after=\par\vspace{\bigskipamount} % Espacio después del cuadro
}


\newcommand{\euler}{e} %Comando para producir letra e de euler
\newenvironment{remark}{\par\vfill\footnotesize % Vertical white space above the remark and smaller font size
\begin{list}{}{
	\leftmargin=80pt % Indentation on the left
	\rightmargin=60pt}\item\ignorespaces % Indentation on the right
\makebox[-2.5pt]{\begin{tikzpicture}[overlay]
		\node[draw=Horizon!60,line width=2.5pt,circle,fill=Horizon!25,font=\sffamily\bfseries,inner sep=4pt,outer sep=0pt] at (-19pt,5pt){\textcolor{Horizon}{Nota}};\end{tikzpicture}} % Blue Nota in a circle
\advance\baselineskip -1pt}{\end{list}\vskip5pt} % Tighter line spacing and white space after remark
\usepackage{graphicx}

\usepackage{titling}

\input{colores.tex}

%%Producir signos de cita textual``Asignatura''

\newcommand{\xmark}{\ding{55}} %%%comando que genera la tachita

\newenvironment{MyColorPar}[1]{%
\leavevmode\color{#1}\ignorespaces%
}{%
}%

%%%%%%%%%%%%%%%%%%% Título de la tarea, Nombre de alumno

\title{ \textcolor{Sun}{\textbf{Ejercicios}}}
\author{\emph{Tarea 3} $\boldsymbol{\mid}$ \textcolor{Tarawera}{\textbf{El problema de Kepler}}}
\date{\today}

%%%%%%%%%%%%%%%%%%% Datos de la Materia

\usepackage{fancyhdr}
\fancypagestyle{plain}{%  the preset of fancyhdr 
\fancyhf{} % clear all header and footer fields
\fancyfoot[R]{
\includegraphics[width=2cm]{LogoFCFMBUAP (1).png}}
\fancyfoot[L]{{\bfseries{\thedate{}}} a las {\bfseries{\mycurrenttime{}}} horas{} (GMT-6, H. Puebla de Zaragoza, Pue)}
\fancyhead[L]{\textbf{Asignatura:} Mecánica Teórica}
\fancyhead[R]{\theauthor}
}
\makeatletter
\def\@maketitle{%
\newpage
\null%
\begin{center}%
\let \footnote \thanks
{\LARGE \@title \par}%
\vskip 1em%
%{\large \@date}%
\end{center}%
}
\makeatother

\usepackage{lipsum}  
\usepackage{cmbright}


%%%%%%%%%%%%%%%%%%%%%%%%%%%%%%%%%%%%%%%%%%%%%%%%%%%%%%%%%%%%%%%%%%%%%%%%%%%%%%%%%%%%%%%%%%%%%%%%%%%%%%%%%%%%%%%%%%%%%%%%%%%%%%%%%%%%%%%%%%%%%%%%%%%%%%%%%%%%%%%%%%%%%%%%%%%%%%%%%%%%%%%%%%%%%%%%%%%%%%

% Por si el color se llama "Apple Green" con espacio en colores.tex:
\colorlet{AppleGreen}{Apple Green}

% Caja AppleGreen con tono ajustable
\newenvironment{AGbox}[1][7]% 7 = tono por defecto (más grande => más claro)
{%
  \begin{tcolorbox}[
    enhanced, breakable,
    colframe=AppleGreen,
    colback=AppleGreen!#1,
    boxrule=1pt,
    borderline west={2mm}{-10mm}{AppleGreen},
    sharp corners=southwest,
    left=14pt,right=14pt,top=10pt,bottom=10pt,
     boxsep=6pt,
  ]%
}{%
  \end{tcolorbox}%
}


\begin{document}



\maketitle

%%%%%%%%%%%%%%%%%%% Datos del profesor, campus, aula e integrantes

\begin{mybox2}
\noindent
\begin{tabular}{@{}p{3.7cm}p{10.4cm}@{}}
\textbf{Profesor} & Dr. Gilberto Silva Ortigoza \\[3pt]
\textbf{Campus / Facultad} & \textcolor{prussianblue}{\textbf{C.U. BUAP}} / \textcolor{trueblue}{\textbf{FCFM}} \\[3pt]
\textbf{Edificio / Salón} & \textcolor{prussianblue}{\textbf{EMA4}} / \textcolor{trueblue}{\textbf{302}} \\[6pt]
\textbf{Integrantes} &
\begin{minipage}[t]{\linewidth}
\raggedright
\textbf{González-Salud}, S.\textsuperscript{1};\,
\textbf{Victoria-García}, A. A.\textsuperscript{2};\,
\textbf{Cruz-Toledo}, F. A.\textsuperscript{3};\,
\textbf{Calderón-Muñoz}, J. C.\textsuperscript{(oyente)4};\,
\textbf{Gómez-González}, F. S.\textsuperscript{5};\,
\textbf{Hernández-Mejía}, M. I.\textsuperscript{6};\,
\textbf{Torres-Luna}, N. A.\textsuperscript{7};\,
\textbf{Martínez-Solorio}, R. J.\textsuperscript{8};\,
\textbf{Arroyo-Gonzalez}, D.\textsuperscript{9};\,
\textbf{Sandoval-Zonotl}, H.\textsuperscript{10};\,
\textbf{Peralta-Rebolledo}, A.\textsuperscript{11};\,
\textbf{Gómez-Alatorre}, A. D.\textsuperscript{12};\,
\textbf{Ballinas-García}, J. A.\textsuperscript{13}
\end{minipage}
\end{tabular}
\end{mybox2}
\vspace{1cm}


\vspace{0.5cm}
\begin{center}
\begin{tikzpicture}
	% Dibuja círculos rellenos en diferentes colores
	\fill[bittersweet] (0,0) circle (0.17cm);
	\fill[Apple Green] (0.6,0) circle (0.17cm);
	\fill[trueblue] (1.2,0) circle (0.17cm);
	\fill[Aguamarina] (1.8,0) circle (0.17cm);
	\fill[Sun] (2.4,0) circle (0.17cm);
	\fill[orange(colorwheel)] (3.0,0) circle (0.17cm);
	\fill[Pesto] (3.6,0) circle (0.17cm);
	\fill[blue-violet] (4.2,0) circle (0.17cm);
	\fill[ticklemepink] (4.8,0) circle (0.17cm);
	\fill[Gallery] (5.4,0) circle (0.17cm);
	\fill[Mantle] (6.0,0) circle (0.17cm);
	\fill[Nevada] (6.6,0) circle (0.17cm);
	\fill[onyx] (7.2,0) circle (0.17cm);
	% Agrega más círculos o modifica los colores según sea necesario
\end{tikzpicture}
\end{center}  \vspace{0.5cm}

\section*{\textcolor{Tarawera}{\textbf{Lista de problemas y soluciones}}}
\addcontentsline{toc}{section}{\textcolor{Tarawera}{\textbf{Problemas: cap 3}}}

\setlength{\columnsep}{2cm}
%\begin{multicols}{2}

\begin{AGbox}[8]
\begin{enumerate}[label={{\textcolor{trueblue}{\textbf{Pr.} \arabic*.1}}}, start=3]
  \item Para el sistema de dos partículas libres demostrar que el momento lineal total, $\boldsymbol{P}$,
  y el momento angular total $\boldsymbol{L}$, están dados por:
  \begin{equation}
  \begin{split}
   \boldsymbol{P}  &= m_{1}\dot{\boldsymbol{r}}_{1} + m_{2}\dot{\boldsymbol{r}}_{2}\\
   \mathbf{L} &= \boldsymbol{r}_{1} \times \boldsymbol{p}_{1} + \boldsymbol{r}_{2} \times \boldsymbol{p}_{2}
  \end{split}
  \end{equation}
  respectivamente. Además, demuestre que son vectores constantes en el tiempo.
\end{enumerate}
\end{AGbox} \vspace{0.5cm}

\textcolor{Cinnabar}{\textbf{Solución:}}

Analizamos los \textbf{grados de libertad} del sistema.  
Al tener dos partículas puntuales ($N=2$) que pueden moverse libremente en el espacio tridimensional ($\mathbb{R}^{3}$) y sin restricciones ($k=0$), el número total de grados de libertad es:
\begin{equation}
n = 3N - k = 3(2) - 0 = 6.
\label{eq:grados-libertad}
\end{equation}
Cada partícula requiere tres coordenadas para describir su posición, por lo que el sistema completo posee seis coordenadas independientes:

\[ \boxed{q_{1} = x_{1}}, \quad \boxed{q_{2}=y_{1}}, \quad \boxed{q_{3}=z_{1}}, \quad \boxed{q_{4}=x_{2}}, \quad \boxed{q_{5}=y_{2}}, \quad \boxed{q_{6}=z_{2}} \]

\vspace{0.3cm}

Denotamos los vectores de posición de las partículas $m_1$ y $m_2$ respecto al origen como:
\begin{equation}
\boldsymbol{r}_1 = (x_1,\, y_1,\, z_1), \qquad
\boldsymbol{r}_2 = (x_2,\, y_2,\, z_2).
\label{eq:posiciones}
\end{equation}
Por tanto, las coordenadas generalizadas y sus derivadas temporales son:
\begin{equation}
\boldsymbol{q} = (x_1,\, y_1,\, z_1,\, x_2,\, y_2,\, z_2), 
\qquad
\dot{\boldsymbol{q}} = (\dot{x}_1,\, \dot{y}_1,\, \dot{z}_1,\, \dot{x}_2,\, \dot{y}_2,\, \dot{z}_2)
\label{eq:coordenadas}
\end{equation}

\vspace{0.3cm}

Al tratarse de \textbf{partículas libres}, la energía potencial $U$ es constante y puede fijarse como $U=0$.  
Así, la Lagrangiana depende únicamente de la energía cinética:
\begin{equation}
\begin{split}
L(\boldsymbol{q}, \dot{\boldsymbol{q}}) 
&= \frac{1}{2}m_1\dot{\boldsymbol{r}}_1^{\,2} + \frac{1}{2}m_2\dot{\boldsymbol{r}}_2^{\,2} \\
&= \frac{1}{2}m_1\left(\dot{x}_1^{\,2} + \dot{y}_1^{\,2} + \dot{z}_1^{\,2}\right)
   + \frac{1}{2}m_2\left(\dot{x}_2^{\,2} + \dot{y}_2^{\,2} + \dot{z}_2^{\,2}\right)
\end{split}
\label{eq:lagrangiana-libre}
\end{equation}

\vspace{0.3cm}

Consideremos ahora las dos partículas libres de masas $m_1$ y $m_2$, con vectores de posición $\boldsymbol{r}_1$ y $\boldsymbol{r}_2$.

\begin{figure}[H]
    \centering
	    \includegraphics[width=0.6\textwidth]{1.png}
    \caption{Representación sencilla del sistema de dos partículas libres en el espacio $\mathbb{R}^{3}$ de masas $m_1$ y $m_2$, 
con sus respectivos vectores de posición $\boldsymbol{r}_1$ y $\boldsymbol{r}_2$ medidos desde el origen. 
El vector relativo $\boldsymbol{r}_2 - \boldsymbol{r}_1$ indica la dirección y separación entre ambas partículas.}
    \label{fig:mi_figura2}
\end{figure}

\vspace{0.3cm}

El \textbf{momento lineal individual} de cada partícula se define como:
\begin{equation}
\boldsymbol{p}_i = m_i \dot{\boldsymbol{r}}_i 
= m_i\big(\dot{x}_i\,\hat{\boldsymbol{i}} + \dot{y}_i\,\hat{\boldsymbol{j}} + \dot{z}_i\,\hat{\boldsymbol{k}}\big),
\qquad i = 1,2.
\label{eq:momento-individual}
\end{equation}


A partir de \eqref{eq:momento-individual}, el \textbf{momento lineal total} del sistema resulta de la suma de los momentos individuales de ambas partículas:

\begin{equation}
\boldsymbol{P} = \boldsymbol{p}_1 + \boldsymbol{p}_2 
= m_1\dot{\boldsymbol{r}}_1 + m_2\dot{\boldsymbol{r}}_2.
\label{eq:momento-total}
\end{equation}

De manera explícita, sustituyendo las componentes cartesianas de cada vector 
\(\dot{\boldsymbol{r}}_i = (\dot{x}_i,\, \dot{y}_i,\, \dot{z}_i)\), 
se obtiene:

\[
\boldsymbol{P} 
= m_1(\dot{x}_1\,\hat{\boldsymbol{i}} + \dot{y}_1\,\hat{\boldsymbol{j}} + \dot{z}_1\,\hat{\boldsymbol{k}}) 
 + m_2(\dot{x}_2\,\hat{\boldsymbol{i}} + \dot{y}_2\,\hat{\boldsymbol{j}} + \dot{z}_2\,\hat{\boldsymbol{k}}).
\]

Agrupando las componentes, el \textbf{momento lineal total} puede escribirse como:

\begin{equation}
\boldsymbol{P} = 
\big(m_1\dot{x}_1 + m_2\dot{x}_2\big)\hat{\boldsymbol{i}} \;+\;
\big(m_1\dot{y}_1 + m_2\dot{y}_2\big)\hat{\boldsymbol{j}} \;+\;
\big(m_1\dot{z}_1 + m_2\dot{z}_2\big)\hat{\boldsymbol{k}}.
\label{eq:momento-total-componentes}
\end{equation}

Por tanto, las \textbf{componentes escalares} del momento lineal total del sistema son:
\[
\boxed{
\begin{aligned}
P_x &= m_1\dot{x}_1 + m_2\dot{x}_2, \\[2pt]
P_y &= m_1\dot{y}_1 + m_2\dot{y}_2, \\[2pt]
P_z &= m_1\dot{z}_1 + m_2\dot{z}_2.
\end{aligned}
}
\]

Cada componente representa la cantidad de movimiento total en su respectiva dirección cartesiana.
En conjunto, el vector \(\boldsymbol{P} = (P_x, P_y, P_z)\) 
describe el impulso total del sistema de dos partículas en el espacio tridimensional.


Para estudiar la evolución temporal de $\boldsymbol{P}$, partimos de \eqref{eq:momento-total} y derivamos con respecto al tiempo:
\begin{equation}
\begin{split}
\dot{\boldsymbol{P}}
&= \frac{d}{dt}\Big(\boldsymbol{p}_1 + \boldsymbol{p}_2\Big)
 = \dot{\boldsymbol{p}}_1 + \dot{\boldsymbol{p}}_2.
\end{split}
\label{eq:derivada-p}
\end{equation}

A partir de la definición del momento lineal individual \eqref{eq:momento-individual},
su derivada temporal componente a componente es:

\begin{equation}
\begin{split}
\dot{\boldsymbol{p}}_i 
&= \frac{d}{dt}(m_i \dot{\boldsymbol{r}}_i) = m_i \ddot{\boldsymbol{r}}_i \\[4pt]
&\Rightarrow 
\begin{cases}
\dot{p}_{ix} = m_i \ddot{x}_i, \\[4pt]
\dot{p}_{iy} = m_i \ddot{y}_i, \\[4pt]
\dot{p}_{iz} = m_i \ddot{z}_i,
\end{cases}
\end{split}
\label{eq:derivada-pi-escalar}
\end{equation}

lo cual, por la \textbf{segunda ley de Newton}, equivale a escribir:

\begin{equation}
\begin{cases}
F_{ix} = m_i \ddot{x}_i, \\[4pt]
F_{iy} = m_i \ddot{y}_i, \\[4pt]
F_{iz} = m_i \ddot{z}_i,
\end{cases}
\qquad \text{o bien, en forma vectorial:} \quad 
\boldsymbol{F}_i = m_i \ddot{\boldsymbol{r}}_i.
\label{eq:newton-escalar}
\end{equation}

De \eqref{eq:derivada-pi-escalar} y \eqref{eq:newton-escalar} se obtiene la relación directa entre la variación temporal del momento y la fuerza aplicada:

\[
\dot{p}_{ix} = F_{ix}, \qquad 
\dot{p}_{iy} = F_{iy}, \qquad 
\dot{p}_{iz} = F_{iz},
\]
oen forma vectorial
\(\dot{\boldsymbol{p}}_i = \boldsymbol{F}_i\).


El \textbf{momento lineal total} del sistema en cada dirección cartesiana se escribe entonces como:

\begin{equation}
\begin{cases}
P_x = p_{1x} + p_{2x} = m_1\dot{x}_1 + m_2\dot{x}_2, \\[4pt]
P_y = p_{1y} + p_{2y} = m_1\dot{y}_1 + m_2\dot{y}_2, \\[4pt]
P_z = p_{1z} + p_{2z} = m_1\dot{z}_1 + m_2\dot{z}_2.
\end{cases}
\label{eq:P-escalar}
\end{equation}

Derivando con respecto al tiempo, se obtiene la variación temporal del momento lineal total:

\begin{equation}
\begin{cases}
\dot{P}_x = \dot{p}_{1x} + \dot{p}_{2x} = F_{1x} + F_{2x}, \\[4pt]
\dot{P}_y = \dot{p}_{1y} + \dot{p}_{2y} = F_{1y} + F_{2y}, \\[4pt]
\dot{P}_z = \dot{p}_{1z} + \dot{p}_{2z} = F_{1z} + F_{2z}.
\end{cases}
\label{eq:variacion-momento-escalar}
\end{equation}

En forma vectorial, esto corresponde a:
\[
\dot{\boldsymbol{P}} = \dot{\boldsymbol{p}}_1 + \dot{\boldsymbol{p}}_2 = \boldsymbol{F}_1 + \boldsymbol{F}_2.
\]

Si el sistema está aislado, sólo existen las fuerzas \textbf{internas} entre las partículas:
\(\boldsymbol{F}_{1\leftarrow2}\) (sobre 1 por 2) y \(\boldsymbol{F}_{2\leftarrow1}\) (sobre 2 por 1).  
Por la \textbf{tercera ley de Newton}:

\[
\boldsymbol{F}_{1\leftarrow2} = -\,\boldsymbol{F}_{2\leftarrow1}.
\]

Entonces, cada componente cumple:

\[
\begin{cases}
F_{1\leftarrow2,\,x} + F_{2\leftarrow1,\,x} = 0, \\[4pt]
F_{1\leftarrow2,\,y} + F_{2\leftarrow1,\,y} = 0, \\[4pt]
F_{1\leftarrow2,\,z} + F_{2\leftarrow1,\,z} = 0,
\end{cases}
\]

lo cual implica directamente que:

\begin{equation}
\dot{P}_x = \dot{P}_y = \dot{P}_z = 0.
\label{eq:conservacion-P-escalar}
\end{equation}

Es decir que si la derivada temporal de una cantidad física es cero, significa que dicha cantidad no cambia con el tiempo, es decir, se \textbf{conserva}.  
Matemáticamente,
\[
\frac{d}{dt}(\text{constante}) = 0,
\]
y físicamente esto se traduce en que el momento lineal total del sistema permanece invariable durante la evolución del movimiento.

Por tanto, el momento lineal total se conserva \emph{en cada dirección independiente del espacio}, lo que expresa que en un sistema aislado ninguna de las componentes \(P_x\), \(P_y\) o \(P_z\) varía en el tiempo:
\[
\boxed{
P_x = \text{cte.}, \quad
P_y = \text{cte.}, \quad
P_z = \text{cte.}
}
\]


\[
\boxed{
\boldsymbol{P} = \text{cte.} \quad \Rightarrow \quad 
(P_x, P_y, P_z) = \text{constantes.}
}
\]


% =========================================================
% Momento angular total: desarrollo con matrices (determinantes)
% =========================================================

Estamos ahora en la parte de demostrar la siguiente ecuación del momento angular total:
\begin{equation}
\boldsymbol{L}=\boldsymbol{r}_1\times\boldsymbol{p}_1+\boldsymbol{r}_2\times\boldsymbol{p}_2
\label{eq:def-L-matrices}
\end{equation}

% =========================================================
% Momento angular total: determinantes, componentes y derivadas
% =========================================================

\textcolor{orange(colorwheel)}{\textbf{Producto cruz mediante determinantes (expansión por cofactores).}} 

Para cada partícula:

\textcolor{Cerulean}{\textbf{Partícula de masa $m_{1}$}}:
\[
\boldsymbol{r}_1\times\boldsymbol{p}_1=
\begin{vmatrix}
\hat{\boldsymbol{i}} & \hat{\boldsymbol{j}} & \hat{\boldsymbol{k}}\\
x_1 & y_1 & z_1\\
p_{1x} & p_{1y} & p_{1z}
\end{vmatrix}
=
(y_1p_{1z}-z_1p_{1y})\,\hat{\boldsymbol{i}}
-(x_1p_{1z}-z_1p_{1x})\,\hat{\boldsymbol{j}}
+(x_1p_{1y}-y_1p_{1x})\,\hat{\boldsymbol{k}}
\]

\textcolor{pakistangreen}{\textbf{Partícula de masa $m_{2}$}}:
\[
\boldsymbol{r}_2\times\boldsymbol{p}_2=
\begin{vmatrix}
\hat{\boldsymbol{i}} & \hat{\boldsymbol{j}} & \hat{\boldsymbol{k}}\\
x_2 & y_2 & z_2\\
p_{2x} & p_{2y} & p_{2z}
\end{vmatrix}
=
(y_2p_{2z}-z_2p_{2y})\,\hat{\boldsymbol{i}}
-(x_2p_{2z}-z_2p_{2x})\,\hat{\boldsymbol{j}}
+(x_2p_{2y}-y_2p_{2x})\,\hat{\boldsymbol{k}}
\]

\textcolor{Apple Green}{\textbf{Componentes escalares de $\boldsymbol L$}} (sumando ambas contribuciones):
\[
\boxed{
\begin{aligned}
L_x &= (y_1p_{1z}-z_1p_{1y})+(y_2p_{2z}-z_2p_{2y})\\[2pt]
L_y &= (z_1p_{1x}-x_1p_{1z})+(z_2p_{2x}-x_2p_{2z})\\[2pt]
L_z &= (x_1p_{1y}-y_1p_{1x})+(x_2p_{2y}-y_2p_{2x})
\end{aligned}
}
\]

% ------------------ Derivadas: paso explícito + cancelaciones justificadas ------------------

\textcolor{vividviolet}{\textbf{Derivadas temporales.}} 
Derivamos cada componente usando la regla del producto:
\[
\begin{cases}
\dot{L}_x = (\dot{y}_1p_{1z}+y_1\dot{p}_{1z})-(\dot{z}_1p_{1y}+z_1\dot{p}_{1y})
           + (\dot{y}_2p_{2z}+y_2\dot{p}_{2z})-(\dot{z}_2p_{2y}+z_2\dot{p}_{2y}), \\[6pt]
\dot{L}_y = (\dot{z}_1p_{1x}+z_1\dot{p}_{1x})-(\dot{x}_1p_{1z}+x_1\dot{p}_{1z})
           + (\dot{z}_2p_{2x}+z_2\dot{p}_{2x})-(\dot{x}_2p_{2z}+x_2\dot{p}_{2z}), \\[6pt]
\dot{L}_z = (\dot{x}_1p_{1y}+x_1\dot{p}_{1y})-(\dot{y}_1p_{1x}+y_1\dot{p}_{1x})
           + (\dot{x}_2p_{2y}+x_2\dot{p}_{2y})-(\dot{y}_2p_{2x}+y_2\dot{p}_{2x})
\end{cases}
\]

Los términos del tipo
\(\dot{y}_i p_{iz} - \dot{z}_i p_{iy}\), \(\dot{z}_i p_{ix} - \dot{x}_i p_{iz}\), 
\(\dot{x}_i p_{iy} - \dot{y}_i p_{ix}\) son \emph{exactamente} las componentes del producto cruz 
\(\dot{\boldsymbol r}_i \times \boldsymbol p_i\):
\[
(\dot{\boldsymbol r}_i \times \boldsymbol p_i)_x=\dot y_i p_{iz}-\dot z_i p_{iy},\quad
(\dot{\boldsymbol r}_i \times \boldsymbol p_i)_y=\dot z_i p_{ix}-\dot x_i p_{iz},\quad
(\dot{\boldsymbol r}_i \times \boldsymbol p_i)_z=\dot x_i p_{iy}-\dot y_i p_{ix}.
\]

Ahora como \(\boldsymbol p_i = m_i \dot{\boldsymbol r}_i\), velocidad y momento son paralelos. 
Por el determinante tenemos:
\[
\dot{\boldsymbol r}_i \times \boldsymbol p_i =
\begin{vmatrix}
\hat{\boldsymbol i} & \hat{\boldsymbol j} & \hat{\boldsymbol k}\\
\dot x_i & \dot y_i & \dot z_i\\
p_{ix} & p_{iy} & p_{iz}
\end{vmatrix}
=
(\dot y_i p_{iz}-\dot z_i p_{iy})\,\hat{\boldsymbol i}
+(\dot z_i p_{ix}-\dot x_i p_{iz})\,\hat{\boldsymbol j}
+(\dot x_i p_{iy}-\dot y_i p_{ix})\,\hat{\boldsymbol k}.
\]

\textbf{Componente $x$}

Sustituyendo \(p_{i\alpha}=m_i \dot r_{i\alpha}\) (con \(\alpha=x,y,z\)), se tiene lo siguiente:
\[
\begin{aligned}
(\dot{\boldsymbol r}_i \times \boldsymbol p_i)_x
&= \dot y_i (m_i \dot z_i) - \dot z_i (m_i \dot y_i) \\[4pt]
&= m_i\,\underbrace{(\dot y_i \dot z_i - \dot z_i \dot y_i)}_{\substack{\text{t\'erminos iguales}\\ \text{i.e. } ab - ba=0}} = 0\\
&= m_i(0) = 0.
\end{aligned}
\]

Así para las expresiones de \(\dot{L}_x,\dot{L}_y,\dot{L}_z\) se cancelan \emph{todos} los términos de la forma 
\((\dot{\boldsymbol r}_i \times \boldsymbol p_i)_\alpha\). 
Sobreviven únicamente los términos con \(\dot{\boldsymbol p}_i\), que, por la segunda ley de Newton, 
son fuerzas: \(\dot{\boldsymbol p}_i=\boldsymbol F_i\). 
Así, reagrupando por componentes,
\[
\boxed{
\begin{aligned}
\dot{L}_x &= (y_1F_{1z}-z_1F_{1y})+(y_2F_{2z}-z_2F_{2y})\\[2pt]
\dot{L}_y &= (z_1F_{1x}-x_1F_{1z})+(z_2F_{2x}-x_2F_{2z})\\[2pt]
\dot{L}_z &= (x_1F_{1y}-y_1F_{1x})+(x_2F_{2y}-y_2F_{2x})
\end{aligned}
}
\]

observamos que cada conjunto de términos tiene la misma estructura que las componentes de un producto vectorial del tipo \(\boldsymbol{r}_i\times\boldsymbol{F}_i\).  
Para verlo de manera explícita, consideremos el desarrollo del producto cruz mediante \textcolor{orange(colorwheel)}{\textbf{expansión por cofactores}}:

\[
\boldsymbol{r}_i\times\boldsymbol{F}_i =
\begin{vmatrix}
\hat{\boldsymbol{i}} & \hat{\boldsymbol{j}} & \hat{\boldsymbol{k}}\\
x_i & y_i & z_i\\
F_{ix} & F_{iy} & F_{iz}
\end{vmatrix}
= (y_iF_{iz}-z_iF_{iy})\,\hat{\boldsymbol{i}}
+ (z_iF_{ix}-x_iF_{iz})\,\hat{\boldsymbol{j}}
+ (x_iF_{iy}-y_iF_{ix})\,\hat{\boldsymbol{k}}.
\]

Cada una de estas componentes coincide exactamente con los términos de las ecuaciones anteriores:  
la primera línea con la componente \(x\), la segunda con la componente \(y\) y la tercera con la componente \(z\).  
Por tanto, podemos reescribir el conjunto de tres ecuaciones escalares en una sola ecuación vectorial.
\[
\boxed{\ \dot{\boldsymbol{L}} = (\boldsymbol{r}_1\times\boldsymbol{F}_1) + (\boldsymbol{r}_2\times\boldsymbol{F}_2)\ }\!,
\]

lo que muestra que la variación temporal del momento angular total es igual a la suma de los \textbf{torques} (momentos de fuerza) ejercidos sobre cada partícula del sistema.


% ------------------ Caso de fuerzas internas centrales ------------------

\textcolor{trueblue}{\textbf{Caso de fuerzas internas centrales.}}

En un sistema aislado, las únicas fuerzas que actúan son \textbf{internas}, es decir, aquellas que las partículas ejercen entre sí.  
Sea \(\boldsymbol{F}_{1\leftarrow 2}\) la fuerza que la partícula 2 ejerce sobre la partícula 1, y \(\boldsymbol{F}_{2\leftarrow 1}=-\boldsymbol{F}_{1\leftarrow 2}\) su correspondiente fuerza de reacción, de acuerdo con la \textbf{tercera ley de Newton}.  

Si además dichas fuerzas son \emph{centrales}, esto es, actúan a lo largo de la línea que une a ambas partículas,
\[
\boldsymbol{F}_{1\leftarrow 2} \parallel \boldsymbol{r} = \boldsymbol{r}_2 - \boldsymbol{r}_1,
\]
entonces el torque interno total del sistema se anula.  
En efecto, partiendo de la relación general
\[
\dot{\boldsymbol{L}} = \boldsymbol{r}_1 \times \boldsymbol{F}_1 + \boldsymbol{r}_2 \times \boldsymbol{F}_2,
\]
y sustituyendo \(\boldsymbol{F}_1 = \boldsymbol{F}_{1\leftarrow 2}\), \(\boldsymbol{F}_2 = \boldsymbol{F}_{2\leftarrow 1} = -\boldsymbol{F}_{1\leftarrow 2}\), se obtiene
\begin{equation}
\begin{split}
\dot{\boldsymbol{L}} 
&= \boldsymbol{r}_1 \times \boldsymbol{F}_{1\leftarrow 2} 
   + \boldsymbol{r}_2 \times \boldsymbol{F}_{2\leftarrow 1} 
   \quad \text{(definición general del torque total)} \\[6pt]
&= \boldsymbol{r}_1 \times \boldsymbol{F}_{1\leftarrow 2} 
   + \boldsymbol{r}_2 \times (-\boldsymbol{F}_{1\leftarrow 2})
   \quad \text{(tercera ley: } \boldsymbol{F}_{2\leftarrow 1}=-\boldsymbol{F}_{1\leftarrow 2}\text{)} \\[6pt]
&= \boldsymbol{r}_1 \times \boldsymbol{F}_{1\leftarrow 2} 
   - \boldsymbol{r}_2 \times \boldsymbol{F}_{1\leftarrow 2}
   \quad \text{(linealidad en el segundo argumento: } \boldsymbol{a}\times(\alpha\,\boldsymbol{b})=\alpha\,(\boldsymbol{a}\times\boldsymbol{b})\text{)} \\[6pt]
&= (\boldsymbol{r}_1 - \boldsymbol{r}_2)\times \boldsymbol{F}_{1\leftarrow 2}
   \quad \text{(distributividad: } \boldsymbol{a}\times\boldsymbol{c} - \boldsymbol{b}\times\boldsymbol{c}
   = (\boldsymbol{a}-\boldsymbol{b})\times\boldsymbol{c}\text{)} \\[6pt]
&= \boldsymbol{r}\times \boldsymbol{F}_{1\leftarrow 2}
   \quad \text{(definición: } \boldsymbol{r}=\boldsymbol{r}_1-\boldsymbol{r}_2\text{)}
\end{split}
\label{eq:torque-interno}
\end{equation}


Como el vector de fuerza \(\boldsymbol{F}_{1\leftarrow 2}\) es paralelo al vector de posición relativa \(\boldsymbol{r}\),
su producto vectorial es nulo:
\[
\boxed{\ \dot{\boldsymbol{L}} = \boldsymbol{0}\ }
\]
Por lo tanto, el \textbf{momento angular total del sistema se conserva}:
\[
\boxed{\ \boldsymbol{L} = \text{constante}\ }
\]

\vspace{0.5cm}

\textcolor{trueblue}{\textbf{Conservación del momento lineal.}}

De manera análoga, para el momento lineal total se cumple
\[
\dot{\boldsymbol{P}} = \boldsymbol{F}_1 + \boldsymbol{F}_2.
\]
Si sólo actúan fuerzas internas, entonces
\(\boldsymbol{F}_{1\leftarrow 2} = -\boldsymbol{F}_{2\leftarrow 1}\), y por tanto
\[
\boxed{\ \dot{\boldsymbol{P}} = \boldsymbol{0} \quad \Rightarrow \quad \boldsymbol{P} = \text{constante}\ }.
\]

\vspace{0.4cm}

\begin{AGbox}[8]
\begin{enumerate}[label={{\textcolor{trueblue}{\textbf{Pr.} \arabic*.2}}}, start=3]
  \item Obtener las ecuaciones: (25), (26), (27) y (29).\\
  
  \textbf{Ecuación (25)}:\\
  \[ \boldsymbol{P} = \frac{\partial L}{\partial \dot{\boldsymbol{R}}} = M\dot{\boldsymbol{R}} = \boldsymbol{p}_{1} + \boldsymbol{p}_{2}   \]
  \textbf{Ecuación (26)}:\\
  \[  \boldsymbol{p} = \frac{\partial L}{\partial \dot{\boldsymbol{r} }} = m\dot{\boldsymbol{r}} = \frac{m_{1}}{M}\boldsymbol{p}_{2} - \frac{m_{2}}{M}\boldsymbol{p}_{1} \]
  \textbf{Ecuación (27)}:\\
  \[  h = \sum\limits_{i=1}^{6}\dot{q_{i}} \frac{\partial L}{ \partial\dot{q}_{i}} - L = \frac{M}{2}\dot{\boldsymbol{R}} \cdot \dot{\boldsymbol{R}} + \frac{m}{2}\dot{\boldsymbol{r}} \cdot \dot{\boldsymbol{r}}\]
  \textbf{Ecuación (29)}:\\
  \[ \boldsymbol{L} = \boldsymbol{R}\times \boldsymbol{P} + \boldsymbol{r}\times \boldsymbol{p}  \]
\end{enumerate}
\end{AGbox}

\textcolor{Cinnabar}{\textbf{Solución}}:

\begin{AGbox}[8]
\begin{enumerate}[label={{\textcolor{trueblue}{\textbf{Pr.} \arabic*.3}}}, start=3]
  \item Para $l$ y $E$ dados, hacer un gráfico de las líneas rectas sobre el plano $xy$ dadas por
  \[
  y = -x\cot \bar{\phi} + \sqrt{ \frac{l^{2}}{2mE} }\, \csc \bar{\phi}
  \]
  para diferentes valores de $\bar{\phi}\in[0,2\pi]$
\end{enumerate}
\end{AGbox}

\textcolor{Cinnabar}{\textbf{Solución:}}

La ecuación dada
\[
y = -x\cot\bar{\phi} + \sqrt{\frac{l^2}{2mE}}\,\csc\bar{\phi}
\]
corresponde al caso libre del \textit{problema de Kepler}, es decir, al movimiento de una partícula con energía total \(E\) y momento angular \(l\) constantes bajo una fuerza central nula. 

Podemos realizar un cambio de variable tomando como la variable $a$ al siguiente término de la ecuación que se nos da en el enunciado, tenemos:

\begin{equation}
a \;\equiv\; \sqrt{\frac{l^2}{2mE}} \;=\; \frac{l}{\sqrt{2mE}}
\label{eq:def-a}
\end{equation}

De la definición \eqref{eq:def-a}, la familia de rectas del enunciado puede escribirse de forma compacta como

\begin{equation}
\begin{split}
{\color{onyx}{y}} 
&= -\,{\color{AppleGreen}{x}}\,{\color{orange}{\cot\bar{\phi}}}
  + {\color{Cerulean}{\sqrt{\tfrac{l^{2}}{2mE}}}}\,
    {\color{Cinnabar}{\csc\bar{\phi}}}\\[4pt]
&= -x{\color{orange}{\cot\bar{\phi}}}
  + {\color{Cerulean}{a}}\,{\color{Cinnabar}{\csc\bar{\phi}}}
  \quad \text{tomando}\quad
  \boxed{{\color{Cerulean}{a=\sqrt{\tfrac{l^{2}}{2mE}}}}}\\[4pt]
&= -x\!\left({\color{orange}{\frac{\cos\bar{\phi}}{\sen\bar{\phi}}}}\right)
  + a\!\left({\color{Cinnabar}{\frac{1}{\sen\bar{\phi}}}}\right)
\end{split}\label{eq:recta-compacta}
\end{equation}

Para visualizar mejor la geometría de la trayectoria, tomamos la ecuación 
en forma compacta~\eqref{eq:recta-compacta} y la multiplicamos por 
\({\color{Cinnabar}{\sen\bar{\phi}}}\). Este paso elimina los denominadores 
y evita las \textbf{indeterminaciones} que aparecen cuando
\(\sen\bar{\phi}=0\) (\(\bar{\phi}=0,\pi,2\pi\)).

Así, al “barrer” el ángulo \(\bar{\phi}\) en el intervalo 
\([0,2\pi]\), obtenemos una descripción continua y simétrica de todas las 
rectas tangentes al círculo \(r=a\), sin asíntotas ni discontinuidades 
en su representación. De esta forma, la ecuación adopta una versión más 
limpia y regular:



\begin{equation}
\begin{split}
{\color{onyx}{y}}
&= -x{\color{orange}{\cot\bar{\phi}}}
   + {\color{Cerulean}{a}}\,{\color{Cinnabar}{\csc\bar{\phi}}}
   \quad \text{(ecuación original de la recta)} \\[6pt]
{\color{Cinnabar}{\sen\bar{\phi}}}\,y
&= -x{\color{orange}{\cos\bar{\phi}}}
   + {\color{Cerulean}{a}}
   \quad \text{(multiplicamos por $\sen\bar{\phi}$ para eliminar denominadores)} \\[6pt]
{\color{AppleGreen}{x}}{\color{orange}{\cos\bar{\phi}}}
+{\color{AppleGreen}{y}}{\color{Cinnabar}{\sen\bar{\phi}}}
&={\color{Cerulean}{a}}
\quad \text{(reordenando términos, obtenemos la forma normal)}
\end{split}
\label{eq:recta-normal}
\end{equation}


Podemos notar que el parámetro 

\[
\boxed{{\color{Cerulean}{a = \frac{l}{\sqrt{2mE}}}}}
\]

permanece constante a lo largo de toda la familia de rectas. 
Esto se debe a que 
\({\color{Cerulean}{\boldsymbol{a}}}\) depende exclusivamente del momento 
angular \({\color{Cerulean}{\boldsymbol{l}}}\) y de la energía total 
\({\color{Cerulean}{\boldsymbol{E}}}\), ambas 
\emph{magnitudes invariantes del movimiento} en un sistema gobernado por 
una fuerza central conservativa. 
En otras palabras, si el sistema conserva \(\boldsymbol{l}\) y \(\boldsymbol{E}\), entonces todas 
las rectas comparten la misma distancia fija 
\({\color{Cerulean}{a}}\) respecto al origen, lo que define la 
envolvente circular de radio \(r = {\color{Cerulean}{a}}\).\\[6pt]

Para ilustrar este comportamiento, podemos asignar 
valores numéricos concretos a las magnitudes físicas. 
Por ejemplo, si tomamos una masa \(\boxed{\boldsymbol{m = 1}}\) (en unidades arbitrarias), una 
energía \(\boxed{\color{Cerulean}\boldsymbol{E = 2}}\) y un momento angular \(\boxed{\color{Cerulean}\boldsymbol{l = 3}}\), obtenemos: 

\[
{\color{Cerulean}{a = \frac{l}{\sqrt{2mE}} = \frac{3}{\sqrt{4}} = 1\text{.}5}}
\]

De este modo, el valor 
\({\color{Cerulean}{a = 1\text{.}5}}\) 
corresponde al radio común a toda la familia de rectas generadas por la ecuación~\eqref{eq:recta-normal}.


Se muestran a continuación algunas rectas particulares de la familia de funciones, 
evaluadas para distintos valores de \(\boldsymbol{\bar{\phi}}\in \left[0,2\pi\right] \), manteniendo fijo el 
valor \({\color{Cerulean}{a = 1\text{.}5}}\): 

\begin{table}[H]
\centering
\caption{Ejemplo de valores para distintas orientaciones \(\bar{\phi}\), 
con \({\color{Cerulean}{a = \tfrac{l}{\sqrt{2mE}} = 1\text{.}5}}\) fijo.}
\label{tab:rectas_phi}
\begin{tabular}{@{}cccc@{}}
\toprule
\(\bar{\phi}\) (rad) & \(\cot \bar{\phi}\) & \(\csc \bar{\phi}\) 
& Ecuación de la recta: \(y = -x \cot\bar{\phi} + 
{\color{Cerulean}{a}}\,\csc\bar{\phi}\) \\ 
\midrule
\(\tfrac{\pi}{6}\)  & \(\sqrt{3}\)         & \(2\)               & \(\color{bittersweet}\boldsymbol{y = -\sqrt{3}\,x + 3\text{.}00}\) \\
\(\tfrac{\pi}{4}\)  & \(1\)                & \(\sqrt{2}\)        & \(\color{Apple Green}\boldsymbol{y = -x + 2\text{.}12}\) \\
\(\tfrac{\pi}{3}\)  & \(1/\sqrt{3}\)       & \(2/\sqrt{3}\)      & \(\color{trueblue}\boldsymbol{y = -0\text{.}577\,x + 1\text{.}73}\) \\
\(\tfrac{2\pi}{3}\) & \(-1/\sqrt{3}\)      & \(2/\sqrt{3}\)      & \(\color{Pesto}\boldsymbol{y = 0\text{.}577\,x + 1\text{.}73}\) \\
\(\tfrac{3\pi}{4}\) & \(-1\)               & \(\sqrt{2}\)        & \(\color{ticklemepink}\boldsymbol{y = x + 2\text{.}12}\) \\
\(\tfrac{5\pi}{6}\) & \(-\sqrt{3}\)        & \(2\)               & \(\color{Sun}\boldsymbol{y = \sqrt{3}\,x + 3\text{.}00}\) \\
\bottomrule
\end{tabular}
\end{table}

Cada fila de la tabla muestra cómo varían la pendiente y la ordenada al 
origen de la recta conforme cambia el ángulo \(\boldsymbol{\bar{\phi}}\).


A continuación se muestra el conjunto de las primeras seis rectas generadas a partir de la ecuación \eqref{eq:recta-normal}, con los valores dados en la tabla \ref{tab:rectas_phi}:

\textbf{Primeras 6 rectas}


\begin{center}
\begin{tikzpicture}
\begin{axis}[
    xlabel={$\boldsymbol{x}$},
    ylabel={$\boldsymbol{y}$},
    xmin=-2.2, xmax=2.2,
    ymin=-2.2, ymax=2.2,
    axis lines=middle,
    grid=none,
    width=15cm,
    height=15cm,
    title={Familia $y=-x\cot\bar{\phi} + a\,\csc\bar{\phi}$ con $a=1.5$},
    legend style={font=\footnotesize, at={(1.03,0.5)}, anchor=west}
]

% Rectas individuales con sus respectivos colores
\addplot[very thick, bittersweet, domain=-2.2:2.2, samples=2] { -x*sqrt(3) + 3.0 };
\addlegendentry{$\bar{\phi}=\tfrac{\pi}{6}$}

\addplot[very thick, Apple Green, domain=-2.2:2.2, samples=2] { -x + 2.12 };
\addlegendentry{$\bar{\phi}=\tfrac{\pi}{4}$}

\addplot[very thick, trueblue, domain=-2.2:2.2, samples=2] { -0.577*x + 1.73 };
\addlegendentry{$\bar{\phi}=\tfrac{\pi}{3}$}

\addplot[very thick, Pesto, domain=-2.2:2.2, samples=2] { 0.577*x + 1.73 };
\addlegendentry{$\bar{\phi}=\tfrac{2\pi}{3}$}

\addplot[very thick, ticklemepink, domain=-2.2:2.2, samples=2] { x + 2.12 };
\addlegendentry{$\bar{\phi}=\tfrac{3\pi}{4}$}

\addplot[very thick, Sun, domain=-2.2:2.2, samples=2] { sqrt(3)*x + 3.0 };
\addlegendentry{$\bar{\phi}=\tfrac{5\pi}{6}$}

\end{axis}
\end{tikzpicture}
\end{center}

\textbf{Primeras 18 rectas}



\begin{center}
\begin{tikzpicture}
\begin{axis}[
    xlabel={$\boldsymbol{x}$},
    ylabel={$\boldsymbol{y}$},
    xmin=-2.2, xmax=2.2,
    ymin=-2.2, ymax=2.2,
    axis lines=middle,
    grid=none,
    width=15cm,
    height=15cm,
    title={Familia $y=-x\cot\bar{\phi} + a\,\csc\bar{\phi}$ con $a=1.5$},
    legend pos=outer north east
]
% ---- parámetro ----
% a = 1.5

% círculo envolvente r=a (opcional: comenta si no lo quieres)
\addplot[very thick, Cerulean, samples=361, domain=0:360]
  ({1.5*cos(x)}, {1.5*sin(x)}) node[pos=0.05,anchor=south west]{};

% rectas para los mismos ángulos de tu tabla
% (PGF usa grados en sin/cos)
\foreach \ph in {30,45,60,120,135,150}{
  \addplot[orange, very thick, domain=-2.2:2.2, samples=2]
    { -x*cos(\ph)/sin(\ph) + 1.5/sin(\ph) };
}

\legend{Envolvente $r=a$}
\end{axis}
\end{tikzpicture}
\end{center}\vspace{0.5cm}


\begin{AGbox}[8]
\begin{enumerate}[label={{\textcolor{trueblue}{\textbf{Pr.} \arabic*.4}}}, start=3]
  \item La ecuación de la cónica, sobre el plano $z=0$, con excentricidad $e$ y uno de sus focos en el origen de coordenadas está dada por:
  
  \[  \frac{1}{r} = C \left[  1 + e\cos\left( \phi - \bar{\phi}  \right)  \right] \]
  
donde $C$ y $\bar{\phi}$ son constantes. Tomar $C=1$, después dar el valor $e=0$ y hacer cuatro gráficos, en diferente color, correspondientes a $\bar{\phi}=0, \frac{ \pi }{2}, \pi, \frac{3\pi}{2}$. Hacer este mismo tipo de gráficos para $e=1$ (parábola), $1/2$ (elipse), 2 (hipérbola).
\end{enumerate}
\end{AGbox}

\textcolor{Cinnabar}{\textbf{Solución}}:

\begin{center}
\begin{tikzpicture}
\begin{axis}[
    xlabel={$\boldsymbol{x}$},
    ylabel={$\boldsymbol{y}$},
    xmin=-2.2, xmax=2.2,
    ymin=-2.2, ymax=2.2,
    axis lines=middle,
    grid=none,
    width=15cm,
    height=15cm,
    title={$r=\tfrac{1}{1+e\cos(\phi-\bar{\phi})}$ con $e=\tfrac12$},
    legend style={font=\footnotesize, at={(1.03,0.5)}, anchor=west},
    unbounded coords=jump,
    samples=1600, domain=-179.9:179.9
]
\def\e{0.5}
\addplot[very thick, vividviolet]
({ (1/(1 + \e*cos(x -   0)))*cos(x) }, { (1/(1 + \e*cos(x -   0)))*sin(x) }); \addlegendentry{$\bar{\phi}=0$}
\addplot[very thick, Cinnabar]
({ (1/(1 + \e*cos(x -  90)))*cos(x) }, { (1/(1 + \e*cos(x -  90)))*sin(x) }); \addlegendentry{$\bar{\phi}=\tfrac{\pi}{2}$}
\addplot[very thick, Apple Green]
({ (1/(1 + \e*cos(x - 180)))*cos(x) }, { (1/(1 + \e*cos(x - 180)))*sin(x) }); \addlegendentry{$\bar{\phi}=\pi$}
\addplot[very thick, Sun]
({ (1/(1 + \e*cos(x - 270)))*cos(x) }, { (1/(1 + \e*cos(x - 270)))*sin(x) }); \addlegendentry{$\bar{\phi}=\tfrac{3\pi}{2}$}
\end{axis}
\end{tikzpicture}
\end{center}

\begin{center}
\begin{tikzpicture}
\begin{axis}[
    xlabel={$\boldsymbol{x}$}, ylabel={$\boldsymbol{y}$},
    xmin=-4, xmax=4, ymin=-4, ymax=4,
    axis lines=middle, grid=none, width=15cm, height=15cm,
    title={$r=\tfrac{1}{1+e\cos(\phi-\bar{\phi})}$ con $e=1$},
    legend style={font=\footnotesize, at={(1.03,0.5)}, anchor=west},
    unbounded coords=jump, samples=2000, domain=-170:170,
    restrict x to domain*=-20:20, restrict y to domain*=-20:20, filter discard warning=false
]
\def\e{1.0}
\addplot[very thick, Blue Sapphire]
({ (1/(1 + \e*cos(x -   0)))*cos(x) }, { (1/(1 + \e*cos(x -   0)))*sin(x) }); \addlegendentry{$\bar{\phi}=0$}
\addplot[very thick, Milano Red]
({ (1/(1 + \e*cos(x -  90)))*cos(x) }, { (1/(1 + \e*cos(x -  90)))*sin(x) }); \addlegendentry{$\bar{\phi}=\tfrac{\pi}{2}$}
\addplot[very thick, Earls Green]
({ (1/(1 + \e*cos(x - 180)))*cos(x) }, { (1/(1 + \e*cos(x - 180)))*sin(x) }); \addlegendentry{$\bar{\phi}=\pi$}
\addplot[very thick, Saffron]
({ (1/(1 + \e*cos(x - 270)))*cos(x) }, { (1/(1 + \e*cos(x - 270)))*sin(x) }); \addlegendentry{$\bar{\phi}=\tfrac{3\pi}{2}$}
\end{axis}
\end{tikzpicture}
\end{center}



\begin{AGbox}[8]
\begin{enumerate}[label={{\textcolor{trueblue}{\textbf{Pr.} \arabic*.5}}}, start=3]
  \item Determinar los focos, vértices y todos los parámetros que caracterizan a las cónicas dadas por las Ecs. (101), (102) y (104).\\
  
    \textbf{Ecuación (101)}:\\
  \[ \bar{x} = \frac{1}{2C} - \left(\frac{C}{2} \right)\bar{y}^{2}  \]
  \textbf{Ecuación (102)}:\\
  \[ \frac{ \left(\bar{x} +ea \right)^{2} }{a^{2}} + \frac{\bar{y}^{2}}{a^{2}(1-e^{2})} =1 \]
    donde:
  \[ a=\frac{1}{C\left(1-e^{2}\right)}  \]
  \textbf{Ecuación (104)}:\\
  \[ \frac{ \left(\bar{x} -eb \right)^{2} }{b^{2}} - \frac{\bar{y}^{2}}{b^{2}(e^{2}-1)} =1 \]
    donde:
    \[ b=\frac{1}{C\left(1-e^{2}\right)}  \]
\end{enumerate}
\end{AGbox}

\textcolor{Cinnabar}{\textbf{Solución}}:

\begin{AGbox}[8]
\begin{enumerate}[label={{\textcolor{trueblue}{\textbf{Pr.} \arabic*.6}}}, start=3]
  \item Suponga que en espacio se tienen dos partículas con vectores de posición $\boldsymbol{r}_{1}$ y $\boldsymbol{r}_{2}$. La partícula uno tiene masa $m_{1}$ y carga $e_{1}$, la partícula dos tiene masa $m_{2}$ y carga $e_{2}$. Usar la ley de gravitación universal de Newton y la ley de Coulomb para determinar la energía potencial del sistema de las dos partículas.
\end{enumerate}
\end{AGbox}

\textcolor{Cinnabar}{\textbf{Solución}}:

La energía potencial gravitatoria:
\[
U_{\text{grav}}(r) = -\frac{G m_1 m_2}{r}.
\]
La energía potencial eléctrica (SI):
\[
U_{\text{Coul}}(r) = \frac{1}{4\pi\varepsilon_0}\frac{e_1 e_2}{r} \equiv k\frac{e_1 e_2}{r}.
\]
Suma:
\[
\boxed{U(r) = -\frac{G m_1 m_2}{r} + k\frac{e_1 e_2}{r} = \frac{1}{r}\bigl(k e_1 e_2 - G m_1 m_2\bigr).}
\]


\begin{AGbox}[8]
\begin{enumerate}[label={{\textcolor{trueblue}{\textbf{Pr.} \arabic*.7}}}, start=3]
  \item En el problema 3.6 suponga que la partícula uno es un protón y la segunda un electrón. Es decir, se tiene el modelo proporcionado por la mecánica clásica del átomo de hidrógeno. Si $F_{Q}$ denota la magnitud de la fuerza de Coulomb y $F_{G}$ denota la magnitud de la fuerza gravitacional entre el protón y el electrón, obtener una expresión para $F_{G}/F_{Q}$. Además, consultar los valores de las constantes que aparecen en este cociente para obtener un valor númerico aproximado.
\end{enumerate}
\end{AGbox}

\textcolor{Cinnabar}{\textbf{Solución}}:


\subsection*{Cálculo numérico}
Usamos valores SI:
\[
\begin{aligned}
G &= 6.67430\times10^{-11}\ \mathrm{m^3\,kg^{-1}s^{-2}},\\
m_p &= 1.6726219\times10^{-27}\ \mathrm{kg},\\
m_e &= 9.10938356\times10^{-31}\ \mathrm{kg},\\
k &= 8.9875517923\times10^9\ \mathrm{N\,m^2/C^2},\\
e &= 1.602176634\times10^{-19}\ \mathrm{C}.
\end{aligned}
\]
Numerador:
\[
G m_p m_e \approx 6.67430\times10^{-11}\cdot 1.6726219\times10^{-27}\cdot 9.10938356\times10^{-31}
\approx 1.01693335\times10^{-67}.
\]
Denominador:
\[
k e^2 \approx 8.9875517923\times10^9\cdot(1.602176634\times10^{-19})^2
\approx 2.30707755\times10^{-28}.
\]
División:
\[
\frac{F_G}{F_Q} \approx \frac{1.01693\times10^{-67}}{2.30708\times10^{-28}}
\approx 4.41\times10^{-40}.
\]
Conclusión: la fuerza gravitatoria es \(\sim 10^{-40}\) veces la fuerza eléctrica entre protón y electrón: insignificante para la física atómica.


\begin{AGbox}[8]
\begin{enumerate}[label={{\textcolor{trueblue}{\textbf{Pr.} \arabic*.8}}}, start=3]
  \item Por el problema 3.7 se concluye que en el átomo de hidrógeno la fuerza gravitacional se puede despreciar con respecto a la fuerza de Coulomb. Así que, suponiendo que la interacción, en el átomo de hidrógeno, está dada por la fuerza de Coulomb, determinar una \textbf{Lagrangiana} que describe la evolución en el espacio del electrón y del protón.
\end{enumerate}
\end{AGbox}

\textcolor{Cinnabar}{\textbf{Solución}}:

Posiciones \(\boldsymbol r_1,\boldsymbol r_2\). Energía cinética:
\[
T = \tfrac12 m_1\dot{\boldsymbol r}_1^2 + \tfrac12 m_2\dot{\boldsymbol r}_2^2.
\]
Potencial Coulomb (protón-electrón, con \(\alpha=k e^2\)):
\[
V(r) = -\frac{\alpha}{r},\qquad r=|\boldsymbol r_2-\boldsymbol r_1|.
\]
Lagrangiana:
\[
\boxed{L = T - V = \tfrac12 m_1\dot{\boldsymbol r}_1^2 + \tfrac12 m_2\dot{\boldsymbol r}_2^2 + \frac{\alpha}{|\boldsymbol r_2-\boldsymbol r_1|}.}
\]


\begin{AGbox}[8]
\begin{enumerate}[label={{\textcolor{trueblue}{\textbf{Pr.} \arabic*.9}}}, start=3]
  \item Obtener las coordenadas cíclicas o ignorables de la \textbf{Lagrangiana} del problema 3.8 y dar su significado físico.
\end{enumerate}
\end{AGbox}

\textcolor{Cinnabar}{\textbf{Solución}}:

Pasar a cilíndricas para cada partícula:
\[
\boldsymbol r_i = (\rho_i\cos\phi_i,\ \rho_i\sin\phi_i,\ z_i),\quad
\dot{\boldsymbol r}_i^2=\dot\rho_i^2+\rho_i^2\dot\phi_i^2+\dot z_i^2.
\]
Lagrangiana:
\[
L=\tfrac12 m_1(\dot\rho_1^2+\rho_1^2\dot\phi_1^2+\dot z_1^2)
+\tfrac12 m_2(\dot\rho_2^2+\rho_2^2\dot\phi_2^2+\dot z_2^2)
- V(r_{12}),
\]
con
\[
r_{12}^2 = \rho_1^2+\rho_2^2-2\rho_1\rho_2\cos(\phi_2-\phi_1)+(z_2-z_1)^2,
\]
de modo que \(V\) depende sólo de \(\Delta\phi=\phi_2-\phi_1\).

Bajo rotación global \(\phi_i\mapsto\phi_i+\varepsilon\) (constante), \(V\) no cambia y \(\dot\phi_i\) tampoco. Por Noether, la cantidad conservada es
\[
J_z = \frac{\partial L}{\partial\dot\phi_1} + \frac{\partial L}{\partial\dot\phi_2}
= m_1\rho_1^2\dot\phi_1 + m_2\rho_2^2\dot\phi_2,
\]
que es el momento angular total alrededor de \(z\).



\begin{AGbox}[8]
\begin{enumerate}[label={{\textcolor{trueblue}{\textbf{Pr.} \arabic*.10}}}, start=3]
  \item Usar las definiciones del vector centro de masas y del vector relativo dadas por:
  \[R = \frac{ m_{1} \boldsymbol{r_{1}} + m_{2}\boldsymbol{r}_{2}  }{m_{1} + m_{2}}, \quad \boldsymbol{r}  = \boldsymbol{r}_{2} - \boldsymbol{r}_{1}     \] 
  para obtener la expresión de la \textbf{Lagrangiana} obtenida en el problema 3.8, en términos de $\boldsymbol{R}$ y $\boldsymbol{r}$. Observe que la \textbf{Lagrangiana} se puede escribir como la suma de dos \textbf{Lagrangianas}, una describe a una partícula libre de masa $M = m_{p} + m_{e}$, y la otra describe una partícula de masa reducida $m=m_{p}m_{e}/(m_{p}+m_{e})\approx m_{e}$ sujeta al potencial de Coulomb. 
\end{enumerate}
\end{AGbox}

\textcolor{Cinnabar}{\textbf{Solución}}:

\[
\boldsymbol R = \frac{m_1\boldsymbol r_1 + m_2\boldsymbol r_2}{M},\qquad
\boldsymbol r = \boldsymbol r_2 - \boldsymbol r_1,\qquad
M=m_1+m_2,\qquad \mu=\frac{m_1m_2}{M}.
\]
Inversas:
\[
\boldsymbol r_1 = \boldsymbol R - \frac{m_2}{M}\boldsymbol r,\qquad
\boldsymbol r_2 = \boldsymbol R + \frac{m_1}{M}\boldsymbol r.
\]

\subsection*{Energía cinética}
Sustituyendo las velocidades:
\[
\dot{\boldsymbol r}_1 = \dot{\boldsymbol R} - \frac{m_2}{M}\dot{\boldsymbol r},\qquad
\dot{\boldsymbol r}_2 = \dot{\boldsymbol R} + \frac{m_1}{M}\dot{\boldsymbol r},
\]
se obtiene (álgebra estándar)
\[
T=\tfrac12 M\dot{\boldsymbol R}^2 + \tfrac12 \mu\dot{\boldsymbol r}^2
\]

\subsection*{Lagrangiana separada}
Como \(V=V(r)\) depende sólo de \(r=|\boldsymbol r|\),
\[
\boxed{L = \tfrac12 M\dot{\boldsymbol R}^2 + \tfrac12 \mu\dot{\boldsymbol r}^2 + \frac{\alpha}{r}.}
\]
Así se separa \(L=L_{\text{CM}}(\boldsymbol R)+L_{\text{rel}}(\boldsymbol r)\)

\begin{AGbox}[8]
\begin{enumerate}[label={{\textcolor{trueblue}{\textbf{Pr.} \arabic*.11}}}, start=3]
  \item Demostrar que el momento angular de la partícula de masa reducida $m=m_{p}m_{e}/(m_{p}+m_{e})\approx m_{e}$, en el problema 3.10, es una constante de movimiento.
\end{enumerate}
\end{AGbox}

\textcolor{Cinnabar}{\textbf{Solución}}:

En el plano (o en coordenadas polares para \(\boldsymbol r\)):
\[
L_{\text{rel}} = \tfrac12\mu(\dot r^2 + r^2\dot\phi^2) + \frac{\alpha}{r}.
\]
\(\phi\) no aparece explícitamente, por tanto es coordenada cíclica. El momento conjugado
\[
p_\phi = \frac{\partial L}{\partial\dot\phi} = \mu r^2\dot\phi \equiv \ell
\]
es constante: \(\ell=\) momento angular de la partícula reducida (conservado).

\begin{AGbox}[8]
\begin{enumerate}[label={{\textcolor{trueblue}{\textbf{Pr.} \arabic*.12}}}, start=3]
  \item Tomar la \textbf{Lagrangiana} que describe la evolución de la partícula de masa reducida en el problema 3.10 y elegir un sistema de referencia de tal forma que la dirección del momento angular orbital de esta partícula coincida con el eje $z$, suponer que la energía de tal partícula es negativa y determine su trayectoria en términos de las condiciones iniciales. Además, hacer un gráfico mostrando las trayectorias para valores dados de la energía y momento angular orbital
\end{enumerate}
\end{AGbox}

\textcolor{Cinnabar}{\textbf{Solución}}:

La energía total relativa:
\[
E = \tfrac12\mu\dot r^2 + \tfrac12\mu r^2\dot\phi^2 - \frac{\alpha}{r}
= \tfrac12\mu\dot r^2 + \frac{\ell^2}{2\mu r^2} - \frac{\alpha}{r}.
\]
Definimos el potencial efectivo
\[
V_{\mathrm{eff}}(r) = \frac{\ell^2}{2\mu r^2} - \frac{\alpha}{r}.
\]

\subsection*{Cambio a variable \(u(\phi)=1/r\) (ecuación de Binet)}
Relación:
\[
\dot r = \frac{dr}{dt} = \frac{dr}{d\phi}\dot\phi = \frac{dr}{d\phi}\frac{\ell}{\mu r^2}.
\]
Con \(u=1/r\), \(dr/d\phi = -u^{-2} du/d\phi\) y
\[
\dot r = -\frac{\ell}{\mu}\frac{du}{d\phi}.
\]
Sustituyendo en \(E\) y manipulando,
\[
\Bigl(\frac{du}{d\phi}\Bigr)^2 + u^2 - \frac{2\mu\alpha}{\ell^2}u = \frac{2\mu E}{\ell^2}.
\]
Derivando esta identidad respecto. a \(\phi\) y reordenando (o usando derivación Newtoniana) se obtiene la ecuación lineal:
\[
\frac{d^2 u}{d\phi^2} + u = \frac{\mu\alpha}{\ell^2}.
\]
Solución general:
\[
u(\phi) = \frac{\mu\alpha}{\ell^2}\Bigl[1 + e\cos(\phi-\phi_0)\Bigr],
\]
o en la forma clásica de la cónica:
\[
\boxed{\displaystyle r(\phi) = \frac{p}{1+e\cos(\phi-\phi_0)},\qquad p=\frac{\ell^2}{\mu\alpha}.}
\]

\subsection*{Relación entre \(e\) y \(E\)}
Sustituyendo la solución en la identidad previa se llega a:
\[
A^2 e^2 = A^2 + \frac{2\mu E}{\ell^2},
\qquad A=\frac{\mu\alpha}{\ell^2}.
\]
Despejando,
\[
\boxed{e^2 = 1 + \frac{2E\ell^2}{\mu\alpha^2}.}
\]
Por tanto:
\begin{itemize}
    \item \(E<0 \Rightarrow e<1\) (elipse),
    \item \(E=0 \Rightarrow e=1\) (parábola),
    \item \(E>0 \Rightarrow e>1\) (hipérbola).
\end{itemize}

\subsection*{Parámetros orbitales}
Semieje mayor \(a\) (para elipses, \(E<0\)):
\[
\frac{E}{\mu} = -\frac{\alpha}{2a} \quad\Rightarrow\quad a = -\frac{\mu\alpha}{2E}.
\]
Relaciones:
\[
p = a(1-e^2),\qquad r_{\min} = a(1-e),\qquad r_{\max}=a(1+e).
\]

\subsection*{Determinación de \(e\) y fase desde condiciones iniciales}
Dado \(r_0\) y \(\dot r_0\) en \(\phi=\phi_{\mathrm{init}}\), con \(u_0=1/r_0\),
\[
\begin{cases}
u_0 = A\bigl(1 + e\cos\delta\bigr),\\[4pt]
\dot r_0 = \dfrac{\ell}{\mu}A e\sin\delta,
\end{cases}
\quad \delta = \phi_{\mathrm{init}} - \phi_0.
\]
De aquí:
\[
e^2 = \Bigl(\frac{u_0}{A}-1\Bigr)^2 + \Bigl(\frac{\mu\dot r_0}{\ell A}\Bigr)^2,
\quad
\delta = \operatorname{atan2}\!\Bigl(\frac{\mu\dot r_0}{\ell A},\,\frac{u_0}{A}-1\Bigr).
\]


\section{ejercicios de funciones especiales}



\end{document}
